\documentclass[UTF8,11pt]{ctexart}
\usepackage[a4paper,margin=24mm]{geometry}
\usepackage{amsmath,amssymb,amsthm,mathtools}
\usepackage[all]{xy}
\usepackage{xurl,hyperref,enumitem}
\hypersetup{hidelinks,breaklinks=true}
\setlength{\emergencystretch}{3em}
\newtheorem*{theorem}{Theorem}
\newtheorem*{lemma}{Lemma}
\newtheorem*{proposition}{Proposition}
\newtheorem*{corollary}{Corollary}
\theoremstyle{definition}
\newtheorem*{definition}{Definition}
\newtheorem*{remark}{Remark}
\DeclareMathOperator{\Hom}{Hom}
\DeclareMathOperator{\Ext}{Ext}
\DeclareMathOperator{\Tor}{Tor}
\DeclareMathOperator{\Spec}{Spec}
\DeclareMathOperator{\supp}{supp}
\DeclareMathOperator{\im}{im}
\DeclareMathOperator{\id}{id}
\DeclareMathOperator{\Tr}{Tr}
\DeclareMathOperator{\rank}{rank}
\newcommand{\R}{\mathbb R}
\newcommand{\C}{\mathbb C}
\newcommand{\Z}{\mathbb Z}
\newcommand{\N}{\mathbb N}
\newcommand{\Q}{\mathbb Q}
\begin{document}
\section*{017\quad 有限表示态射的可构造像定理}
\subsection*{来源与收录范围}
The Stacks Project Authors, \emph{The Stacks Project}. Chapters 10 and 29，主目标为 Theorem 29.23.3；其核心仿射证明 Theorem 10.29.10 与 Lemmas 10.29.7--10.29.9 一并收录，只计一个 Chevalley 定理问题。
\par\url{https://stacks.math.columbia.edu/tag/054K}
\par\url{https://stacks.math.columbia.edu/tag/00FE}
\par\url{https://stacks.math.columbia.edu/tag/00F5}
\par 使用 2026-09-28 实际访问的在线正文；未取得网页构建的固定版本号。以网页数学 TeX 整理，不把本摘录称为整章下载原件。
\subsection*{作者命题与人类证明}

\begin{lemma}[10.29.7]
Let $R$ be a ring. The map $\Spec(R[x])\to\Spec(R)$ is open, and the image of any standard open is a quasi-compact open.
\end{lemma}
\begin{proof}
It suffices to show that the image of a standard open $D(f)$, $f\in R[x]$ is quasi-compact open. The image of $D(f)$ is the image of $\Spec(R[x]_f)\to\Spec(R)$. Let $\mathfrak p\subset R$ be a prime ideal. Let $\overline f$ be the image of $f$ in $\kappa(\mathfrak p)[x]$. Recall, see Lemma 10.18.6, that $\mathfrak p$ is in the image if and only if $R[x]_f\otimes_R\kappa(\mathfrak p)=\kappa(\mathfrak p)[x]_{\overline f}$ is not the zero ring. This is exactly the condition that $f$ does not map to zero in $\kappa(\mathfrak p)[x]$, in other words, that some coefficient of $f$ is not in $\mathfrak p$. Hence we see: if $f=a_dx^d+\ldots+a_0$, then the image of $D(f)$ is $D(a_d)\cup\ldots\cup D(a_0)$.
\end{proof}
\begin{lemma}[10.29.8]
Let $R\to A$ be a ring homomorphism. Assume $A\cong R^{\oplus n}$ as an $R$-module. Let $f\in A$. The multiplication map $m_f:A\to A$ is $R$-linear and hence has a characteristic polynomial $P(T)=T^n+r_{n-1}T^{n-1}+\ldots+r_0\in R[T]$. For any prime $\mathfrak p\in\Spec(R)$, $f$ acts nilpotently on $A\otimes_R\kappa(\mathfrak p)$ if and only if $\mathfrak p\in V(r_0,\ldots,r_{n-1})$.
\end{lemma}
\begin{proof}
This follows quite easily once we prove that the characteristic polynomial $\bar P(T)\in\kappa(\mathfrak p)[T]$ of the multiplication map $m_{\bar f}:A\otimes_R\kappa(\mathfrak p)\to A\otimes_R\kappa(\mathfrak p)$ which multiplies elements of $A\otimes_R\kappa(\mathfrak p)$ by $\bar f$, the image of $f$ viewed in $\kappa(\mathfrak p)$, is just the image of $P(T)$ in $\kappa(\mathfrak p)[T]$.

Let $(a_{ij})$ be the matrix of the map $m_f$ with entries in $R$, using a basis $e_1,\ldots,e_n$ of $A$ as an $R$-module. Then, $A\otimes_R\kappa(\mathfrak p)\cong(R\otimes_R\kappa(\mathfrak p))^{\oplus n}=\kappa(\mathfrak p)^n$, which is an $n$-dimensional vector space over $\kappa(\mathfrak p)$ with basis $e_1\otimes1,\ldots,e_n\otimes1$. The image $\bar f=f\otimes1$, and so the multiplication map $m_{\bar f}$ has matrix $(a_{ij}\otimes1)$. Thus, the characteristic polynomial is precisely the image of $P(T)$.

From linear algebra, we know that a linear transformation acts nilpotently on an $n$-dimensional vector space if and only if the characteristic polynomial is $T^n$ (since the characteristic polynomial divides some power of the minimal polynomial). Hence, $f$ acts nilpotently on $A\otimes_R\kappa(\mathfrak p)$ if and only if $\bar P(T)=T^n$. This occurs if and only if $r_i\in\mathfrak p$ for all $0\leq i\leq n-1$, that is when $\mathfrak p\in V(r_0,\ldots,r_{n-1})$.
\end{proof}
\begin{lemma}[10.29.9]
Let $R$ be a ring. Let $f,g\in R[x]$ be polynomials. Assume the leading coefficient of $g$ is a unit of $R$. There exists elements $r_i\in R$, $i=1\ldots,n$ such that the image of $D(f)\cap V(g)$ in $\Spec(R)$ is $\bigcup_{i=1,\ldots,n}D(r_i)$.
\end{lemma}
\begin{proof}
Write $g=ux^d+a_{d-1}x^{d-1}+\ldots+a_0$, where $d$ is the degree of $g$, and hence $u\in R^*$. Consider the ring $A=R[x]/(g)$. It is, as an $R$-module, finite free with basis the images of $1,x,\ldots,x^{d-1}$. Consider multiplication by (the image of) $f$ on $A$. This is an $R$-module map. Hence we can let $P(T)\in R[T]$ be the characteristic polynomial of this map. Write $P(T)=T^d+r_{d-1}T^{d-1}+\ldots+r_0$.

We claim that $r_0,\ldots,r_{d-1}$ have the desired property. We will use below the property of characteristic polynomials that
\[
\mathfrak p\in V(r_0,\ldots,r_{d-1})\quad\Longleftrightarrow\quad
\text{multiplication by $f$ is nilpotent on }A\otimes_R\kappa(\mathfrak p).
\]
This was proved in Lemma 10.29.8.

Suppose $\mathfrak q\in D(f)\cap V(g)$, and let $\mathfrak p=\mathfrak q\cap R$. Then there is a nonzero map $A\otimes_R\kappa(\mathfrak p)\to\kappa(\mathfrak q)$ which is compatible with multiplication by $f$. And $f$ acts as a unit on $\kappa(\mathfrak q)$. Thus we conclude $\mathfrak p\notin V(r_0,\ldots,r_{d-1})$.

On the other hand, suppose that $r_i\notin\mathfrak p$ for some prime $\mathfrak p$ of $R$ and some $0\leq i\leq d-1$. Then multiplication by $f$ is not nilpotent on the algebra $A\otimes_R\kappa(\mathfrak p)$. Hence there exists a prime ideal $\overline{\mathfrak q}\subset A\otimes_R\kappa(\mathfrak p)$ not containing the image of $f$. The inverse image of $\overline{\mathfrak q}$ in $R[x]$ is an element of $D(f)\cap V(g)$ mapping to $\mathfrak p$.
\end{proof}
\begin{theorem}[10.29.10, Chevalley's Theorem, Tag 00FE]
Suppose that $R\to S$ is of finite presentation. The image of a constructible subset of $\Spec(S)$ in $\Spec(R)$ is constructible.
\end{theorem}
\begin{proof}
Write $S=R[x_1,\ldots,x_n]/(f_1,\ldots,f_m)$. We may factor $R\to S$ as $R\to R[x_1]\to R[x_1,x_2]\to\ldots\to R[x_1,\ldots,x_{n-1}]\to S$. Hence we may assume that $S=R[x]/(f_1,\ldots,f_m)$. In this case we factor the map as $R\to R[x]\to S$, and by Lemma 10.29.6 we reduce to the case $S=R[x]$.

By Lemma 10.29.1 suffices to show that if $T=(\bigcup_{i=1\ldots n}D(f_i))\cap V(g_1,\ldots,g_m)$ for $f_i,g_j\in R[x]$ then the image in $\Spec(R)$ is constructible. Since finite unions of constructible sets are constructible, it suffices to deal with the case $n=1$, i.e., when $T=D(f)\cap V(g_1,\ldots,g_m)$.

Note that if $c\in R$, then we have
\[\Spec(R)=V(c)\amalg D(c)=\Spec(R/(c))\amalg\Spec(R_c),\]
and correspondingly $\Spec(R[x])=V(c)\amalg D(c)=\Spec(R/(c)[x])\amalg\Spec(R_c[x])$. The intersection of $T=D(f)\cap V(g_1,\ldots,g_m)$ with each part still has the same shape, with $f,g_i$ replaced by their images in $R/(c)[x]$, respectively $R_c[x]$. Note that the image of $T$ in $\Spec(R)$ is the union of the image of $T\cap V(c)$ and $T\cap D(c)$. Using Lemmas 10.29.5 and 10.29.6 it suffices to prove the images of both parts are constructible in $\Spec(R/(c))$, respectively $\Spec(R_c)$.

Let us assume we have $T=D(f)\cap V(g_1,\ldots,g_m)$ as above, with $\deg(g_1)\leq\deg(g_2)\leq\ldots\leq\deg(g_m)$. We are going to use induction on $m$, and on the degrees of the $g_i$. Let $d_1=\deg(g_1)$, i.e., $g_1=cx^{d_1}+l.o.t$ with $c\in R$ not zero. Cutting $R$ up into the pieces $R/(c)$ and $R_c$ we either lower the degree of $g_1$ (and this is covered by induction) or we reduce to the case where $c$ is invertible.

If $c$ is invertible, and $m>1$, then write $g_2=c'x^{d_2}+l.o.t$. In this case consider $g_2'=g_2-(c'/c)x^{d_2-d_1}g_1$. Since the ideals $(g_1,g_2,\ldots,g_m)$ and $(g_1,g_2',g_3,\ldots,g_m)$ are equal we see that $T=D(f)\cap V(g_1,g_2',g_3\ldots,g_m)$. But here the degree of $g_2'$ is strictly less than the degree of $g_2$ and hence this case is covered by induction.

The bases case for the induction above are the cases (a) $T=D(f)\cap V(g)$ where the leading coefficient of $g$ is invertible, and (b) $T=D(f)$. These two cases are dealt with in Lemmas 10.29.9 and 10.29.7.
\end{proof}
\begin{theorem}[29.23.3, Chevalley's Theorem, Tag 054K]
Let $f:X\to Y$ be a morphism of schemes. Assume $f$ is quasi-compact and locally of finite presentation. Then the image of every locally constructible subset is locally constructible.
\end{theorem}
\begin{proof}
Let $E\subset X$ be locally constructible. We have to show that $f(E)$ is locally constructible too. We will show that $f(E)\cap V$ is constructible for any affine open $V\subset Y$. Thus we reduce to the case where $Y$ is affine. In this case $X$ is quasi-compact. Hence we can write $X=U_1\cup\ldots\cup U_n$ with each $U_i$ affine open in $X$. If $E\subset X$ is locally constructible, then each $E\cap U_i$ is constructible, see Properties, Lemma 28.2.1. Hence, since $f(E)=\bigcup f(E\cap U_i)$ and since finite unions of constructible sets are constructible, this reduces us to the case where $X$ is affine. In this case the result is Algebra, Theorem 10.29.10.
\end{proof}

\subsection*{整理说明（非作者正文）}
记号据源书解释：$D(f)$ 是 $f$ 非零的标准开集，$V(I)$ 是闭集，$\kappa(\mathfrak p)$ 是剩余域。在环谱上，可构造集是有限个 $D(f)\cap V(g_1,\ldots,g_m)$ 的并（10.29.3）；一般概形使用局部可构造概念。这里没有额外假设 $R$ Noetherian。

标准先修是上述可构造集基本表示、主开局部化及有限生成理想取商对可构造集的性质（10.29.1--10.29.6）、剩余域纤维非空判据。关键步骤——特征多项式的系数控制纤维中幂零性、首项系数分层及多项式消元归纳——均由正文承担而未当作先修略去。概形版与仿射版不是两道题。

保留作者原语言、论证次序及数学内容；仅调整独立文档的排版和外部引用显示。原有章节号均指源书，不是本文件编号。网页评论不属于作者证明，未收录。
\subsection*{可修改的简短标注（非作者正文）}
$c$：有限表示态射下集合像的可构造性

$a$：剩余域纤维；有限自由代数；乘法算子特征多项式；幂零性；局部化与取商分层；多项式消元；仿射覆盖

\texttt{cross\_theory\_status = yes}。将像的存在性条件转换为有限自由代数中乘法算子的非幂零性，再用特征多项式系数给出开集条件；一般情况通过系数分层与降次归纳处理，并返回概形局部可构造性。位置：10.29.8--10.29.10 与 29.23.3。

全部标注均为非穷尽、可修改初稿；未标注不表示未使用。
\subsection*{许可与修改记录}
原数学正文为 The Stacks Project Authors 的作品，按 GNU Free Documentation License 1.2 或后续版本复用；无不变章节、封面文字或封底文字。许可全文在 \texttt{sources/stacks/GFDL-1.2.txt}。本整理沿用原许可。2026-09-28：助手摘录、独立排版并加入来源、范围说明与初稿标注；不暗示原作者认可本次整理。
\end{document}
